% three-column role page: #1 colour, #2 "Bạn thấy", #3 "Việc của bạn bây giờ", #4 "Sắp có"
\newcommand\rolecols[4]{%
\begin{tikzpicture}[x=1mm,y=1mm]
\draw[#1,line width=0.8pt,fill=#1!6,rounded corners=3pt] (0,0) rectangle (42,-62);
\node[anchor=north west,font=\scriptsize\bfseries,text=#1!85!black] at (2,-2) {Bạn thấy};
\node[anchor=north west,font=\fsz{6.2}{7.8},text width=38mm] at (2,-8) {#2};
\draw[ink!30,line width=0.8pt,fill=white,rounded corners=3pt] (45,0) rectangle (105,-62);
\node[anchor=north west,font=\scriptsize\bfseries] at (47,-2) {Việc của bạn bây giờ};
\node[anchor=north west,font=\fsz{6.2}{7.8},text width=56mm] at (47,-8) {#3};
\draw[warn!60,line width=0.8pt,dashed,fill=soft,rounded corners=3pt] (108,0) rectangle (144,-62);
\node[anchor=north west,font=\scriptsize\bfseries,text=warn!80!black] at (110,-2) {Sắp có trên trang web};
\node[anchor=north west,font=\fsz{6}{7.6},text width=32mm,text=ink!85] at (110,-8) {#4};
\end{tikzpicture}}

%======================================================== PB
\begin{frame}[label=pb]{\rtag{PB}\ Phản biện}
\rolecols{PB}{%
• Chỉ \textbf{bài được giao cho bạn} trong kỳ đang mở.\\[1mm]
• Đề bài đã biên tập; \textbf{lời giải} khi PT mở (sau khi bạn tự giải).\\[1mm]
• Thảo luận của bài đó; nhận xét của phản biện khác hiện \textbf{ẩn danh} (\emph{Phản biện 1, 2…}).\\[3mm]
\textbf{Không thấy}: tên tác giả, tên tệp, nguồn, các bài khác — để chấm \textbf{ẩn danh}.\\[3mm]
Hết kỳ, bài rời khỏi danh sách của bạn.}{%
\textbf{1.}\ Nhận email mời, mở \textbf{đường dẫn đăng nhập}. Mỗi bài có chip \emph{hạn …}.\\[1mm]
\textbf{2.}\ Mở từng bài: \textbf{tự giải trước}; lời giải hiện khi PT mở cho kỳ.\\[1mm]
\textbf{3.}\ Điền \textbf{Phiếu phản biện} dưới lời giải, bấm \emph{Lưu phiếu}:\\[0.5mm]
\hspace*{3mm}• Mức đề nghị: \textbf{A} hay \textbf{B}?\\
\hspace*{3mm}• Đề nghị: \textbf{chọn} / \textbf{sửa rồi chọn} / \textbf{không chọn}.\\
\hspace*{3mm}• Nhận xét: đề rõ, đủ điều kiện chưa? Lời giải đúng, đủ chưa? Giống bài đã biết?\\[1mm]
\textbf{4.}\ Bấm \textbf{Đánh dấu đã xong} (phiếu bị khoá; bỏ đánh dấu để sửa).\\[1mm]
\textbf{5.}\ Trao đổi với người khác ở \emph{Thảo luận} của bài.\\[2mm]
\textcolor{ink!65}{Chưa xong thì có email nhắc 3 và 1 ngày trước hạn. Công thức viết như \LaTeX\ (mục 7).}}{%
• Sửa, xoá nhận xét đã gửi ở Thảo luận.}
\end{frame}

%======================================================== NCB
\begin{frame}[label=ncb]{\rtag{NCB}\ Người chuẩn bị bài}
\rolecols{NCB}{%
• \textbf{Mọi bài}, tên tác giả (không có liên hệ).\\[1mm]
• \textbf{Nguồn \& chỉnh sửa} của từng bài (trang sau).\\[1mm]
• Chấm \textbf{cần kiểm tra n}, \textbf{xung đột mở n} trên danh sách.\\[1mm]
• Thảo luận của mọi bài.}{%
\textbf{Nguyên tắc}\\[0.5mm]
• \textbf{Không sửa lén}: mỗi sửa đổi là một dòng ghi — vị trí, trước, sau, lý do, trạng thái. Văn bản đề, lời giải không chứa ghi chú.\\
• \textbf{Bản gốc của tác giả giữ nguyên}; chỉ sửa bản biên tập.\\
• \textbf{Nghi ngờ thì ghi, không tự sửa}: thành mục \emph{cần kiểm tra}; khi cần, hỏi tác giả qua Văn phòng / PT.\\
• Xung đột ghi theo loại, mỗi cái một dòng.\\[2mm]
\textbf{Bây giờ}\\[0.5mm]
• Lọc danh sách, xử lý các bài có \emph{cần kiểm tra}.\\
• Đối chiếu bản gốc khi \emph{Đối chiếu nguyên bản} còn „chưa"; báo \emph{Cảnh báo hiển thị}.\\
• \textbf{Sửa}: \emph{Sửa nhỏ}, hoặc \emph{Sửa nội dung toán} (vị trí, lý do → Sửa đổi).\\
• Kiểm tra, xung đột: thêm, đóng trong Nguồn \& chỉnh sửa.\\
• \textbf{Hình}: mục Hình (TikZ / tên ảnh); nút \emph{Hình} để dựng SVG.}{%
• Dựng hình tự động (không cần chạy lệnh ở máy).}
\end{frame}

%======================================================== NGUỒN & CHỈNH SỬA
\begin{frame}[label=nguon]{\rtag{NCB}\ \rtag{PT}\ \rtag{TBT}\ Đọc mục „Nguồn \& chỉnh sửa"}
\begin{tikzpicture}[x=1mm,y=1mm]
\draw[win] (0,0) rectangle (70,-63);
\node[anchor=west,font=\fsz{5.4}{6.4}\bfseries] at (2,-3.5) {{\symf ▾} Nguồn \& chỉnh sửa};
\begin{scope}[every node/.style={anchor=north west,font=\fsz{4.8}{5.8},text width=65mm}]
\node at (2,-7) {\textbf{Nguồn}\\ 2026-09 · Thư mục hồ sơ 11 · bai-de-xuat.pdf · 2026-09-14 · email};
\node at (2,-15) {\textbf{Đối chiếu nguyên bản}\\ chưa — chép từ ảnh/PDF};
\node at (2,-23) {\textbf{Sửa đổi}\\ \textbf{lời giải, kết luận}: $x=\frac{4}{3}\to x=\frac{2}{3}$ (bước trước tính ra $\frac{2}{3}$, kết luận chép nhầm; đã sửa ở bản biên tập — chờ tác giả xác nhận)};
\node at (2,-35) {\textbf{Cần kiểm tra}\\ • Công thức chép từ ảnh, cần đối chiếu [mở]\\ • Cần lập luận vì sao mẫu số khác $0$ [mở]};
\node at (2,-47) {\textbf{Xung đột}\\ \textbf{mức} — tác giả đề xuất B, phản biện đề nghị A $\Rightarrow$ (chưa) [mở]};
\node at (2,-56) {\textbf{Nhật ký chuyển đổi}\\ đổi \textbackslash over thành \textbackslash frac (chỉ cách gõ)};
\end{scope}
\begin{scope}[every node/.style={anchor=north west,font=\fsz{5.8}{7.2},text width=68mm}]
\node at (76,-1) {\textbf{Nguồn} — hồ sơ gốc ở đâu, nhận ngày nào, qua kênh nào, các bản trung gian, bài đã qua những vòng nào.};
\node at (76,-10) {\textbf{Đối chiếu nguyên bản} — văn bản trong hệ thống đã được so với tệp gốc của tác giả chưa, ai so, kết quả.};
\node at (76,-19) {\textbf{Sửa đổi} — mọi chỗ bản biên tập khác bản gốc: \emph{vị trí: trước → sau (lý do; trạng thái)}. Trạng thái cho biết tác giả đã đồng ý chưa.};
\node at (76,-31) {\textbf{Cần kiểm tra} — câu hỏi còn mở: nghi lỗi, thiếu hình, chép từ ảnh… [mở] / [xong].};
\node at (76,-40) {\textbf{Xung đột} — mỗi loại một dòng, kèm cách giải quyết:};
\end{scope}
\foreach \t/\d [count=\i] in {số hiệu/{cùng bài, số khác nhau ở các danh sách},mức/{A hay B},bản chép khác nhau/{các bản sao không khớp},trùng bài/{giống bài khác / bài đã biết},đề và lời giải không khớp/{},tác giả/{sáng tác hay sưu tầm, tên in},trạng thái/{danh sách ghi „chính thức" nhưng chưa in}}{
  \node[anchor=west,font=\fsz{5.4}{6.4}] at (79,{-45.5-(\i-1)*2.9}) {• \textbf{\t}\ \textcolor{ink!60}{\d}};}
\end{tikzpicture}
\end{frame}

%======================================================== PT
\begin{frame}[label=pt]{\rtag{PT}\ Phụ trách chuyên mục}
\rolecols{PT}{%
• \textbf{Mọi bài}, tên và \textbf{liên hệ} tác giả.\\[1mm]
• Nguồn \& chỉnh sửa, xung đột, mục cần kiểm tra.\\[1mm]
• Thảo luận của mọi bài.}{%
• \textbf{Việc còn mở}: lọc theo trạng thái, xem chấm \emph{cần kiểm tra}, \emph{xung đột mở}.\\[1mm]
• Xung đột trong phạm vi PT (số hiệu, bản chép, trùng bài…); \emph{mức}, \emph{tác giả}: chuyển TBT.\\[1mm]
• Liên hệ tác giả khi NCB cần hỏi.\\[1mm]
• Người mới: báo Quản trị email và vai trò. Sửa đề, lời giải như NCB.\\[1mm]
• \textbf{Đổi trạng thái} (Mới → SL → SL-OK → PL; loại: Không SL, SL-Fail — ghi lý do).\\[1mm]
• \textbf{Kỳ phản biện} (nút đầu trang): mở kỳ, đặt hạn, giao bài, \emph{Gửi thư mời}, xem phiếu và tiến độ, đóng kỳ.\\[1mm]
• \textbf{Bảng chọn bài}: lập bảng cho số báo, chọn bài cho từng vị trí (4 B, 6 A), \emph{Gửi TBT duyệt}.\\[1mm]
• \textbf{Khoá kỳ…} (bảng đã duyệt): số in, thứ tự, lưu ý; bấm hai lần — bài thành PL.}{%
• Dựng hình tự động.}
\end{frame}

%======================================================== TBT
\begin{frame}[label=tbt]{\rtag{TBT}\ Tổng biên tập}
\rolecols{TBT}{%
• \textbf{Mọi bài}, tên và \textbf{liên hệ} tác giả.\\[1mm]
• Nguồn \& chỉnh sửa, xung đột, mục cần kiểm tra.\\[1mm]
• Thảo luận của mọi bài, gồm nhận xét của phản biện.\\[1mm]
• Phiếu phản biện (cuối trang bài) và tiến độ: nút \emph{Kỳ phản biện}.}{%
\textbf{Quyết định thuộc TBT}\\[0.5mm]
• \textbf{Mức} A / B khi tác giả, NCB và phản biện khác nhau.\\
• \textbf{Tác giả}: sáng tác hay sưu tầm, cách ghi tên khi đăng.\\
• Bài nghi \textbf{trùng} bài đã biết: giữ, sửa hay bỏ.\\
• \textbf{Duyệt cuối} danh sách đăng.\\[2mm]
\textbf{Bây giờ}\\[0.5mm]
• Lọc danh sách; mở bài có \emph{xung đột mở} để xem loại xung đột.\\
• Xung đột „chờ TBT": \textbf{Ghi quyết định…} trong Nguồn \& chỉnh sửa (cách giải quyết, trạng thái).\\
• Đổi trạng thái bài; trạng thái sửa đổi khi tác giả trả lời.\\
• \textbf{Bảng chọn bài}: \emph{Duyệt} (bài chọn thành SL-OK, nhận mức của vị trí) hoặc \emph{Trả lại…} kèm lý do; \emph{Mở lại} để thay bài.\\
• \textbf{Khoá kỳ…} sau khi duyệt: đánh số in (P…), khoá — không mở lại được nữa.\\
• Muốn xem hệ thống như phản biện thấy: nhờ Quản trị dùng \emph{xem như vai trò}.}{%
• Thông báo email gọn: bài nào chờ TBT quyết định.}
\end{frame}

%======================================================== VP
\begin{frame}[label=vp]{\rtag{VP}\ Văn phòng}
\rolecols{VP}{%
• \textbf{Mọi bài}, tên và \textbf{liên hệ} tác giả.\\[1mm]
• Nguồn của từng bài: thư mục hồ sơ, tệp gốc, ngày nhận, kênh.}{%
• Nhận hồ sơ như hiện nay (email, Drive), mỗi hồ sơ một thư mục.\\[1mm]
• \textbf{Thêm bài} (nút đầu trang): tác giả, tệp .tex, ảnh — hệ thống cấp mã bài.\\[1mm]
• Tác giả hỏi về bài: tìm theo mã bài trong danh sách, xem trạng thái.\\[2mm]
\textbf{Thông tin cá nhân}\\[0.5mm]
• Điện thoại, email, địa chỉ của tác giả chỉ nằm trong mục liên hệ (chỉ VP, PT, TBT, Quản trị xem), \textbf{không} chép vào đề, lời giải, nhận xét.\\
• Thông tin của học sinh dưới 18 tuổi: ghi tối thiểu.}{%
• Thư báo nhận bài cho tác giả (tuỳ chọn).\\[1mm]
• Tải hồ sơ .docx của tác giả thẳng vào Thêm bài.}
\end{frame}

%======================================================== BTK
\begin{frame}[label=btk]{\rtag{BTK}\ BTK — chế bản}
\rolecols{BTK}{%
• \textbf{Mọi bài} (đọc).\\[1mm]
• Nguồn \& chỉnh sửa (để biết bài nào còn chờ tác giả xác nhận).}{%
• Xem trước các bài trạng thái \textbf{SL-OK} (chờ đăng): lọc \emph{Mọi trạng thái} → SL-OK.\\[1mm]
• Góp ý về chế bản (hình, công thức dài, cách ngắt dòng) trong Thảo luận của bài.\\[1mm]
• Báo NCB nếu công thức trên trang hiện khác bản in mong muốn.\\[1mm]
• Số đã khoá: nút \emph{Tải tệp .tex} trên trang \textbf{Bảng chọn bài}.}{%
• \textbf{Mỗi kỳ}: một tệp \LaTeX\ \textbf{chỉ có đề bài} để chế bản, theo mẫu cột \emph{Đề ra kỳ này} — số in đã đánh, hình TikZ, dòng tác giả; kèm thư mục \texttt{pic/} và bản PDF xem thử.}
\end{frame}

%======================================================== QT
\begin{frame}[label=qt]{\rtag{QT}\ Quản trị}
\begin{tikzpicture}[x=1mm,y=1mm]
\node[anchor=north west,font=\scriptsize\bfseries] at (0,0) {Thêm người dùng — Sheet „Pi ĐRKN — dữ liệu", tab \texttt{Users}};
\draw[ink!30,fill=white] (0,-5) rectangle (96,-23);
\fill[soft] (0,-5) rectangle (96,-9);
\foreach \x/\h in {0/email,34/ten,52/vai\_tro,74/hoat\_dong,86/ghi\_chu}{\node[anchor=west,font=\fsz{5.4}{6.4}\bfseries\ttfamily] at (\x+1,-7) {\h};}
\foreach \y/\a/\b/\c/\d in {-11/ten.nguoi@example.com/Họ Tên/PB/TRUE,-15/nguoi.khac@example.com/Họ Tên/{NCB, PT}/TRUE,-19/da.nghi@example.com/Họ Tên/PB/\textcolor{warn}{FALSE}}{
  \node[anchor=west,font=\fsz{5.4}{6.4}\ttfamily] at (1,\y) {\a};\node[anchor=west,font=\fsz{5.4}{6.4}] at (35,\y) {\b};
  \node[anchor=west,font=\fsz{5.4}{6.4}] at (53,\y) {\c};\node[anchor=west,font=\fsz{5.4}{6.4}\ttfamily] at (75,\y) {\d};}
\foreach \x in {34,52,74,86}{\draw[ink!20] (\x,-5) -- (\x,-23);}
\node[anchor=north west,font=\fsz{6}{7.6},text width=96mm] at (0,-26) {%
• \textbf{Một người một dòng.} Nhiều vai trò cách nhau dấu phẩy: \texttt{NCB, PT}. Vai trò hợp lệ: TBT, PT, NCB, PB, VP, BTK, Quản trị.\\
• \textbf{Tạm ngưng}: đổi \texttt{hoat\_dong} thành FALSE — có hiệu lực ngay, kể cả khi người đó đang mở trang. Không cần xoá dòng.\\
• Sau khi thêm: gửi cho người đó \textbf{đường dẫn đăng nhập} (không gửi địa chỉ trang hệ thống).\\
• Phân công phản biện: trang \emph{Kỳ phản biện}. Nhắc hạn tự động: chạy \texttt{installReminderTrigger} một lần.\\
• Nhập bài từ tệp: \texttt{importLatest} nhập mọi tệp \texttt{pi-drkn-import…} chưa nhập.\\
• \textbf{Không sửa tay} các tab khác (Problems, Corrections…): sửa qua hệ thống để có lịch sử.};
\draw[QT,line width=0.8pt,fill=QT!5,rounded corners=3pt] (102,0) rectangle (144,-62);
\node[anchor=north west,font=\scriptsize\bfseries,text=QT] at (104,-2) {Xem như vai trò};
\node[anchor=north west,font=\fsz{6}{7.6},text width=38mm] at (104,-8) {Ô \emph{Vai trò thật} ở đầu danh sách → chọn \emph{Xem như PB} (hay vai trò khác) để thấy đúng những gì vai trò đó thấy. Chọn lại \emph{Vai trò thật} để trở về. Chỉ đổi cách xem, không đổi quyền thật.};
\node[anchor=north west,font=\scriptsize\bfseries,text=QT] at (104,-34) {Việc kỹ thuật};
\node[anchor=north west,font=\fsz{6}{7.6},text width=38mm] at (104,-40) {Nhập hàng loạt, bản vá dữ liệu, kiểm thử, bài luyện tập: chạy trong trình soạn thảo Apps Script bằng tài khoản chủ — xem tài liệu kỹ thuật (\texttt{docs/}). Tài khoản chủ bật \textbf{xác minh 2 bước}; không gửi mật khẩu qua email.};
\end{tikzpicture}
\end{frame}

%======================================================== SỰ CỐ
\begin{frame}[label=suco]{Khi gặp sự cố}
\vskip1mm
{\fsz{6.2}{7.6}
\setlength{\tabcolsep}{5pt}\renewcommand{\arraystretch}{1.5}
\begin{tabular}{@{}p{50mm}p{42mm}p{47mm}@{}}
\textbf{Bạn thấy} & \textbf{Vì sao} & \textbf{Làm gì}\\\hline
„Địa chỉ … chưa có vai trò trong hệ thống" & email chưa được thêm, hoặc đang dùng nhầm tài khoản & xem dòng \emph{Xin chào}; nếu đúng tài khoản, báo PT\\
„Phiên vào hệ thống đã hết hạn hoặc không hợp lệ" & tải lại trang sau 10 phút, hoặc mở dấu trang cũ & bấm \textbf{Đăng nhập lại}\\
Dòng đỏ „Phiên làm việc đã hết hạn — hãy vào lại…" & tab đã mở quá 6 giờ & vào lại từ đường dẫn đăng nhập\\
„Không có quyền xem bài này" & bài không giao cho bạn, hoặc kỳ đã đóng & báo PT nếu bạn cần bài đó\\
Trang trắng & trình duyệt tải lại khung trang & đóng tab, vào lại từ đường dẫn đăng nhập\\
Công thức hiện nguyên \texttt{\$…\$}, hoặc dòng \emph{Cảnh báo hiển thị} & lệnh \LaTeX\ hệ thống chưa hiển thị được & ghi vào Thảo luận của bài để NCB xử lý\\
„Bài vừa được người khác sửa (phiên bản n) — bản của bạn CHƯA được lưu" & hai người sửa cùng lúc & bản bạn gõ vẫn trong ô; so với bản mới hiện bên dưới, gộp rồi bấm \textbf{Lưu} lần nữa\\
Dải xanh „This application was created by a Google Apps Script user" & thông báo chuẩn của Google & bấm \no\ để đóng\\\hline
\end{tabular}}
\end{frame}

%======================================================== SẮP CÓ
\begin{frame}[label=sapco]{Đang làm tiếp — hướng dẫn sẽ được cập nhật}
\begin{tikzpicture}[x=1mm,y=1mm]
\foreach \t/\who/\d [count=\i] in {
  {Thêm bài từ tệp .docx}/{VP, NCB, PT}/{tải hồ sơ Word của tác giả thẳng vào Thêm bài},
  {Dựng hình tự động}/{NCB, PT}/{TikZ thành SVG trong một kho GitHub riêng tư, không chạy lệnh ở máy}}{
  \node[badge=NCB,minimum size=5.5mm,font=\fsz{7}{8}\bfseries] at (3,{-4-(\i-1)*9.5}) {\i};
  \node[anchor=west,font=\fsz{7}{8.4}\bfseries] at (8,{-2.6-(\i-1)*9.5}) {\t};
  \node[anchor=west,font=\fsz{6.2}{7.4},text=ink!70] at (8,{-6-(\i-1)*9.5}) {\d};
  \node[anchor=west,font=\fsz{6}{7},text=ink!55] at (86,{-4-(\i-1)*9.5}) {cho: \who};}
\node[anchor=north west,font=\fsz{6.2}{7.6},text=ink!70,text width=140mm] at (0,-58) {Trong kỳ thử 10/2026, phần nào chưa có trên trang web vẫn làm như cũ; mọi góp ý về cách dùng, xin ghi vào Thảo luận của bài luyện tập hoặc báo PT.};
\end{tikzpicture}
\end{frame}
